% GENERATED FILE. Edit canonical lessons, then run npm run generate:books.
% canonical-source-hash: fnv1a64:b422c88e00a46683
% canonical-lessons: PA-C214-ki-return, PA-C214-ki-polar, PA-C214-polar-review, PA-C214-answers, PA-C214-payoff

\chapter{Are You a Traveller?}
\label{ch:pa-a2-polar}
\hlchaptermodality{214}

\begin{chapteropening}
\textbf{By the end of this chapter:} \emph{I can ask a yes-or-no question with ki at the front and answer it with han ji or nahin ji.}

Ask a yes-or-no question with ki, tell it from ki 'what', and answer yes or no politely.
\end{chapteropening}

\section[\texorpdfstring{\pa{ਤੁਹਾਡਾ} \pa{ਨਾਂ} \pa{ਕੀ} \pa{ਹੈ}? … (tuhāḍā nāṁ kī …)}{ਤੁਹਾਡਾ ਨਾਂ ਕੀ ਹੈ? … (tuhāḍā nāṁ kī …)}]{Return to a familiar question}

\label{lesson:PA-C214-ki-return}



\begin{quote}
\emph{Hear:} \emph{tuhāḍā nāṁ kī hai?} Say \textbf{\pa{ਤੁਹਾਡਾ} \pa{ਨਾਂ} \pa{ਕੀ} \pa{ਹੈ}?}, \textquotedblleft{}what is your name?\textquotedblright{}
\end{quote}

\begin{grammarlens}[title={Grammar Lens}]
In that question \textbf{\pa{ਕੀ}} means \textquotedblleft{}what\textquotedblright{}, and it sits where the answer will go. You also know \textbf{\pa{ਯਾਤਰੀ}}, \textquotedblleft{}a traveller\textquotedblright{}, and \textbf{\pa{ਤੁਸੀਂ} … \pa{ਹੋ}}, \textquotedblleft{}you are\textquotedblright{}.
\end{grammarlens}

\subsection*{Your turn}
\begin{itemize}
\raggedright
  \item \emph{Hear:} \emph{what is your name?} Say \textbf{\pa{ਤੁਹਾਡਾ} \pa{ਨਾਂ} \pa{ਕੀ} \pa{ਹੈ}?}
  \item \emph{Hear:} \emph{a traveller} Say \textbf{\pa{ਯਾਤਰੀ}}.
  \item \emph{Hear:} \emph{how are you?} Say \textbf{\pa{ਤੁਸੀਂ} \pa{ਕਿਵੇਂ} \pa{ਹੋ}?}
\end{itemize}

\begin{tcolorbox}[breakable,colback=teal!4,colframe=teal!35!black,title={Before you move on}]
Say \emph{what}, then \emph{a traveller}.
\end{tcolorbox}
\section[\texorpdfstring{\pa{ਕੀ} \pa{ਤੁਸੀਂ} \pa{ਯਾਤਰੀ} \pa{ਹੋ}? (kī tusīṁ yātrī ho?)}{ਕੀ ਤੁਸੀਂ ਯਾਤਰੀ ਹੋ? (kī tusīṁ yātrī ho?)}]{\pa{ਕੀ} \pa{ਤੁਸੀਂ} \pa{ਯਾਤਰੀ} \pa{ਹੋ}? — are you a traveller?}

\label{lesson:PA-C214-ki-polar}



\begin{quote}
Say \emph{what is your name?} Then \emph{you are a traveller}: \textbf{\pa{ਤੁਸੀਂ} \pa{ਯਾਤਰੀ} \pa{ਹੋ}}.
\end{quote}

\begin{grammarlens}[title={Grammar Lens: \pa{ਕੀ} at the front}]
Put \textbf{\pa{ਕੀ}} in front of a statement and it becomes a yes-or-no question. The voice rises at the end:

\begin{quote}
\textbf{\pa{ਤੁਸੀਂ} \pa{ਯਾਤਰੀ} \pa{ਹੋ।}} — \emph{tusīṁ yātrī ho.} — You are a traveller.
\end{quote}

\begin{quote}
\textbf{\pa{ਕੀ} \pa{ਤੁਸੀਂ} \pa{ਯਾਤਰੀ} \pa{ਹੋ}?} — \emph{kī tusīṁ yātrī ho?} — Are you a traveller?
\end{quote}

\begin{quote}
\textbf{\pa{ਕੀ} \pa{ਬਾਗ਼} \pa{ਖੁੱਲ੍ਹਾ} \pa{ਹੈ}?} — \emph{kī bāgh khullhā hai?} — Is the garden open?
\end{quote}

Same word, two jobs. In \textbf{\pa{ਤੁਹਾਡਾ} \pa{ਨਾਂ} \pa{ਕੀ} \pa{ਹੈ}?} it stands where the answer goes, and means \textquotedblleft{}what\textquotedblright{}. At the very front it asks only yes or no.
\end{grammarlens}

\subsection*{Your turn}
\begin{itemize}
\raggedright
  \item \emph{Hear:} \emph{is he a traveller?} Say \textbf{\pa{ਕੀ} \pa{ਉਹ} \pa{ਯਾਤਰੀ} \pa{ਹੈ}?}
  \item \emph{Hear:} \emph{is the garden closed?} Say \textbf{\pa{ਕੀ} \pa{ਬਾਗ਼} \pa{ਬੰਦ} \pa{ਹੈ}?}
  \item \emph{Read it:} \textbf{\pa{ਕੀ} \pa{ਤੁਸੀਂ} \pa{ਯਾਤਰੀ} \pa{ਹੋ}?} Point to \textbf{\pa{ਯ}} in \textbf{\pa{ਯਾਤਰੀ}} and say \emph{ya}.
  \item \emph{Write it:} copy \textbf{\pa{ਕੀ} \pa{ਤੁਸੀਂ} \pa{ਯਾਤਰੀ} \pa{ਹੋ}?} once
\end{itemize}

\begin{tcolorbox}[breakable,colback=teal!4,colframe=teal!35!black,title={Before you move on}]
\textbf{\pa{ਕੀ}} at the front: what does it ask? (Yes or no.) In the middle? (What.)
\end{tcolorbox}
\section[\texorpdfstring{\pa{ਕੀ} \pa{ਤੁਸੀਂ} \pa{ਯਾਤਰੀ} \pa{ਹੋ}? (kī tusīṁ yātrī ho?)}{ਕੀ ਤੁਸੀਂ ਯਾਤਰੀ ਹੋ? (kī tusīṁ yātrī ho?)}]{Ask yes or no again}

\label{lesson:PA-C214-polar-review}



\begin{quote}
\emph{Hear:} \emph{are you a traveller?} Say \textbf{\pa{ਕੀ} \pa{ਤੁਸੀਂ} \pa{ਯਾਤਰੀ} \pa{ਹੋ}?}
\end{quote}

\subsection*{Your turn}
\begin{itemize}
\raggedright
  \item \emph{Hear:} \emph{is she a traveller?} Say \textbf{\pa{ਕੀ} \pa{ਉਹ} \pa{ਯਾਤਰੀ} \pa{ਹੈ}?}
  \item \emph{Hear:} \emph{is the garden open?} Say \textbf{\pa{ਕੀ} \pa{ਬਾਗ਼} \pa{ਖੁੱਲ੍ਹਾ} \pa{ਹੈ}?}
  \item \emph{Read it:} \textbf{\pa{ਯਾਤਰੀ}} Point to \textbf{\pa{ਯ}}, say \emph{ya}, then say the whole word.
  \item \emph{Write it:} cover the model and write \textbf{\pa{ਕੀ} \pa{ਉਹ} \pa{ਯਾਤਰੀ} \pa{ਹੈ}?} once
\end{itemize}

\begin{tcolorbox}[breakable,colback=teal!4,colframe=teal!35!black,title={Before you move on}]
Ask someone whether they are a traveller.
\end{tcolorbox}
\section[\texorpdfstring{\pa{ਹਾਂ} \pa{ਜੀ} · \pa{ਨਹੀਂ} \pa{ਜੀ} (hāṁ jī · nahīṁ jī)}{ਹਾਂ ਜੀ · ਨਹੀਂ ਜੀ (hāṁ jī · nahīṁ jī)}]{\pa{ਹਾਂ} \pa{ਜੀ}, \pa{ਨਹੀਂ} \pa{ਜੀ} — yes, no}

\label{lesson:PA-C214-answers}



\begin{quote}
Say \emph{are you a traveller?} Then \emph{yes}, \emph{no}, and the polite \textbf{\pa{ਜੀ}}.
\end{quote}

\begin{grammarlens}[title={Grammar Lens: answering}]
Answer with \textbf{\pa{ਹਾਂ} \pa{ਜੀ}} or \textbf{\pa{ਨਹੀਂ} \pa{ਜੀ}}, then say the whole sentence back. A no answer puts \textbf{\pa{ਨਹੀਂ}} before the verb, as you already do:

\begin{quote}
\textbf{\pa{ਹਾਂ} \pa{ਜੀ}, \pa{ਮੈਂ} \pa{ਯਾਤਰੀ} \pa{ਹਾਂ।}} — \emph{hāṁ jī, maiṁ yātrī hāṁ.} — Yes, I am a traveller.
\end{quote}

\begin{quote}
\textbf{\pa{ਨਹੀਂ} \pa{ਜੀ}, \pa{ਮੈਂ} \pa{ਯਾਤਰੀ} \pa{ਨਹੀਂ} \pa{ਹਾਂ।}} — \emph{nahīṁ jī, maiṁ yātrī nahīṁ hāṁ.} — No, I am not a traveller.
\end{quote}
\end{grammarlens}

\subsection*{Your turn}
\begin{itemize}
\raggedright
  \item \emph{Hear:} \emph{is the garden open? Yes.} Say \textbf{\pa{ਹਾਂ} \pa{ਜੀ}, \pa{ਬਾਗ਼} \pa{ਖੁੱਲ੍ਹਾ} \pa{ਹੈ।}}
  \item \emph{Hear:} \emph{is the garden closed? No.} Say \textbf{\pa{ਨਹੀਂ} \pa{ਜੀ}, \pa{ਬਾਗ਼} \pa{ਬੰਦ} \pa{ਨਹੀਂ} \pa{ਹੈ।}}
  \item \emph{Read it:} \textbf{\pa{ਨਹੀਂ} \pa{ਜੀ}, \pa{ਮੈਂ} \pa{ਯਾਤਰੀ} \pa{ਨਹੀਂ} \pa{ਹਾਂ।}} Give its meaning.
  \item \emph{Write it:} copy \textbf{\pa{ਹਾਂ} \pa{ਜੀ}} once
\end{itemize}

\begin{tcolorbox}[breakable,colback=teal!4,colframe=teal!35!black,title={Before you move on}]
\emph{Are you a traveller?} Answer it truthfully, in a whole sentence.
\end{tcolorbox}
\section[\texorpdfstring{\pa{ਕੀ} \pa{ਤੁਸੀਂ} \pa{ਯਾਤਰੀ} \pa{ਹੋ}? (kī tusīṁ yātrī ho?)}{ਕੀ ਤੁਸੀਂ ਯਾਤਰੀ ਹੋ? (kī tusīṁ yātrī ho?)}]{Yes or no: a four-skill payoff}

\label{lesson:PA-C214-payoff}



\begin{quote}
Ask \emph{are you a traveller?} Then answer it with \textbf{\pa{ਹਾਂ} \pa{ਜੀ}}.
\end{quote}

\subsection*{Your turn: four skills}
\begin{itemize}
\raggedright
  \item \emph{Hear:} \emph{is he a traveller?} Say \textbf{\pa{ਕੀ} \pa{ਉਹ} \pa{ਯਾਤਰੀ} \pa{ਹੈ}?}
  \item \emph{Hear:} \emph{no, he is not a traveller} Say \textbf{\pa{ਨਹੀਂ} \pa{ਜੀ}, \pa{ਉਹ} \pa{ਯਾਤਰੀ} \pa{ਨਹੀਂ} \pa{ਹੈ।}}
  \item \emph{Read it:} \textbf{\pa{ਕੀ} \pa{ਬਾਗ਼} \pa{ਖੁੱਲ੍ਹਾ} \pa{ਹੈ}?} Give its meaning, then point to \textbf{\pa{ਯ}} in \textbf{\pa{ਯਾਤਰੀ}} once more.
  \item \emph{Write it:} cover the line and write \textbf{\pa{ਕੀ} \pa{ਤੁਸੀਂ} \pa{ਯਾਤਰੀ} \pa{ਹੋ}?}
\end{itemize}

\begin{tcolorbox}[breakable,colback=teal!4,colframe=teal!35!black,title={Before you move on}]
\textbf{\pa{ਕੀ}} asks; \textbf{\pa{ਹਾਂ} \pa{ਜੀ}} or \textbf{\pa{ਨਹੀਂ} \pa{ਜੀ}} answers. Ask and answer once more.
\end{tcolorbox}
